\sezione{3}{La rete del lab basic-ipv4}{topologia e indirizzi}

\subsection{Chi c'è nella rete}
Due router (\texttt{r1}, \texttt{r2}) e tre host (\texttt{pc1}, \texttt{pc2}, \texttt{pc3}) collegati da tre LAN (collision domain \texttt{A}, \texttt{B}, \texttt{C}). Nel laboratorio imparerai a:
\begin{itemize}
  \item assegnare un indirizzo IPv4 e una netmask a un'interfaccia;
  \item impostare il default gateway di un host;
  \item impostare la tabella di routing di un router;
  \item usare \texttt{ping} e \texttt{traceroute} e osservare il comportamento di ARP.
\end{itemize}

\subsection{La topologia completa}
\begin{center}
\topologia{1}{1}{1}
\end{center}
\didascalia{In giallo i collision domain (LAN), in verde le reti IP, in rosso l'ultimo ottetto dell'indirizzo di ogni interfaccia, in viola i MAC address. Per esempio \texttt{r1} ha \texttt{eth0} = \texttt{195.11.14.1} (MAC \texttt{00:..:a1}) su A ed \texttt{eth1} = \texttt{100.0.0.9} (MAC \texttt{00:..:b1}) su B.}

\begin{center}
\rowcolors{2}{chiaro}{white}
\begin{tabular}{l l c l l l}
\rowcolor{navy}\thd{Dispositivo} & \thd{Interfaccia} & \thd{LAN} & \thd{Indirizzo IPv4} & \thd{MAC} & \thd{Default gateway}\\
pc1 & eth0 & A & 195.11.14.5/24 & 00:00:00:00:00:01 & 195.11.14.1 (r1)\\
pc2 & eth0 & C & 200.1.1.7/24 & 00:00:00:00:00:02 & 200.1.1.1 (r2)\\
pc3 & eth0 & C & 200.1.1.3/24 & 00:00:00:00:00:03 & 200.1.1.1 (r2)\\
r1 & eth0 & A & 195.11.14.1/24 & 00:00:00:00:00:a1 & —\\
r1 & eth1 & B & 100.0.0.9/30 & 00:00:00:00:00:b1 & —\\
r2 & eth0 & C & 200.1.1.1/24 & 00:00:00:00:00:c1 & —\\
r2 & eth1 & B & 100.0.0.10/30 & 00:00:00:00:00:b2 & —\\
\end{tabular}
\end{center}

\subsection{Capire le maschere: /24 e /30}
\begin{tecnica}[Notazione CIDR]
Il numero dopo la barra dice quanti bit (da sinistra) identificano la \textbf{rete}; i restanti identificano l'\textbf{host}.
\begin{itemize}
  \item \texttt{/24} = maschera \texttt{255.255.255.0}: 8 bit di host, $2^8=256$ indirizzi. Il primo (\texttt{.0}) è l'indirizzo di rete, l'ultimo (\texttt{.255}) il broadcast: restano \textbf{254} indirizzi utilizzabili.
  \item \texttt{/30} = maschera \texttt{255.255.255.252}: 2 bit di host, 4 indirizzi. Restano \textbf{2} indirizzi utilizzabili: è la rete perfetta per un collegamento punto-punto tra due router.
\end{itemize}
\end{tecnica}
\begin{esempio}[La LAN B: 100.0.0.8/30]
\begin{center}
\begin{tikzpicture}[c/.style={minimum width=2.6cm,minimum height=.8cm,font=\ttfamily\small,draw=white,line width=1.5pt}]
  \node[c,fill=grigio!30] at (0,0) {100.0.0.8};
  \node[c,fill=teal!25] at (2.6,0) {100.0.0.9};
  \node[c,fill=teal!25] at (5.2,0) {100.0.0.10};
  \node[c,fill=grigio!30] at (7.8,0) {100.0.0.11};
  \foreach \x/\t in {0/indirizzo di rete,2.6/r1 eth1,5.2/r2 eth1,7.8/broadcast}{\node[font=\headfont\scriptsize,text=grigio] at (\x,-.65) {\t};}
\end{tikzpicture}
\end{center}
In binario l'ultimo ottetto è \texttt{000010\textbf{00}}…\texttt{000010\textbf{11}}: i primi 6 bit sono fissi (parte di rete), gli ultimi 2 variano. Ecco perché nella tabella di \texttt{r1} comparirà \texttt{100.0.0.11} come indirizzo di broadcast della LAN B.
\end{esempio}

\subsection{Chi sa cosa: la logica del routing}
\begin{intuizione}
Ogni dispositivo conosce «da solo» solo le reti a cui è attaccato direttamente. Per tutto il resto serve un'indicazione esplicita.
\begin{itemize}
  \item \textbf{Gli host} sanno raggiungere la propria LAN; per tutto il resto usano il \textbf{default gateway} («se non sai dove mandarlo, dallo al router»).
  \item \textbf{r1} conosce A e B (direttamente connesse); per C ha una \textbf{rotta statica} via r2.
  \item \textbf{r2} conosce B e C; per A ha una rotta statica via r1.
\end{itemize}
Con queste informazioni ogni host può raggiungere ogni altro host, in entrambe le direzioni.
\end{intuizione}
\begin{center}
\rowcolors{2}{chiaro}{white}
\begin{tabular}{l c c c}
\rowcolor{navy}\thd{} & \thd{LAN A (195.11.14.0/24)} & \thd{LAN B (100.0.0.8/30)} & \thd{LAN C (200.1.1.0/24)}\\
pc1 & diretta & default via r1 & default via r1\\
pc2, pc3 & default via r2 & default via r2 & diretta\\
r1 & diretta (eth0) & diretta (eth1) & \textbf{statica} via 100.0.0.10\\
r2 & \textbf{statica} via 100.0.0.9 & diretta (eth1) & diretta (eth0)\\
\end{tabular}
\end{center}

\begin{attenzione}[Il ritorno conta quanto l'andata]
Un ping funziona solo se esiste la rotta \emph{anche per la risposta}. Se togli la rotta statica da r2, pc1 riesce ancora a mandare la richiesta a pc2 (grazie a r1), ma la risposta di pc2 si ferma su r2, che non sa dove sia \texttt{195.11.14.0/24}. È uno degli errori più frequenti nei laboratori di routing.
\end{attenzione}

\begin{appunti}{Riepilogo del capitolo}
3 LAN: A = 195.11.14.0/24 (pc1 .5, r1 .1); B = 100.0.0.8/30 (r1 .9, r2 .10, broadcast .11); C = 200.1.1.0/24 (pc2 .7, pc3 .3, r2 .1). Host: rete diretta + default gateway. Router: due reti dirette + una rotta statica verso la LAN remota. Le rotte servono in entrambe le direzioni.
\end{appunti}
